\documentclass[11pt]{article}
\usepackage[T1]{fontenc}
\usepackage{lmodern,amsmath,amssymb,amsthm,geometry,hyperref}
\geometry{margin=1in}
\hypersetup{hidelinks}
\newtheorem{theorem}{Theorem}
\newtheorem{lemma}{Lemma}
\newcommand{\RR}{\mathbb R}
\title{Formal verification of two arithmetic-function ratio inequalities}
\author{A verification note on arguments posted by Shivam Patel}
\date{25 October 2026}
\begin{document}
\maketitle
\begin{abstract}
We verify Conjectures 3 and 9 of S. I. Dimitrov's preprint
\emph{Lower bounds on expressions depending on the functions
$\varphi(n)$, $\psi(n)$ and $\sigma(n)$, III}, arXiv:2606.12484v1.
Both statements hold for every integer $n\geq2$ and $k\geq1$.
Shivam Patel posted complete general arguments for these two statements
on MathDB on 20 August 2026. This note credits those arguments and gives
a derivative-free proof with a complete Lean certificate, including the
arithmetic-function bridge. No novelty or equality-classification claim
is made.
\end{abstract}

\section{Statements and arithmetic functions}
For $n\geq1$, write $n=\prod_{p\mid n}p^{e_p}$ with $e_p\geq1$.
Let $\varphi$ be Euler's totient, let $\sigma(n)=\sum_{d\mid n}d$, and define
\[
 \psi(n)=\prod_{p\mid n}p^{e_p-1}(p+1).
\]
Empty products have value one. This is precisely the definition of
Dedekind psi in Section 1 of Dimitrov's preprint. It also satisfies
\[
 \psi(n)=n\prod_{p\mid n}(1+1/p),
\]
since each factor on the right is
$p^{e_p}(1+1/p)=p^{e_p-1}(p+1)$.

\begin{theorem}[Conjecture 3; MathDB 375256]\label{thm:simple}
For all integers $n\geq2$ and $k\geq1$,
\begin{align*}
&\frac{\varphi(n)^k}{\psi(n)^k+\sigma(n)^k}
+\frac{\psi(n)^k}{\varphi(n)^k+\sigma(n)^k}
+\frac{\sigma(n)^k}{\varphi(n)^k+\psi(n)^k}\\
&\hspace{10mm}\geq
\frac{4(n+1)^{2k}+(n^2-1)^k+(n-1)^{2k}}
 {2((n+1)^{2k}+(n^2-1)^k)}.
\end{align*}
\end{theorem}

\begin{theorem}[Conjecture 9; MathDB 375257]\label{thm:mixed}
For all integers $n\geq2$ and $k\geq1$, put
\[
x=\varphi(n)^k,\quad y=\psi(n)^k,\quad z=\sigma(n)^k.
\]
Then
\begin{align*}
&\frac{x(y+z)}{y(x+z)+z(x+y)}
 +\frac{y(x+z)}{x(y+z)+z(x+y)}
 +\frac{z(x+y)}{x(y+z)+y(x+z)}\\
&\hspace{10mm}\geq
\frac{5(n-1)^{2k}+5(n^2-1)^k+2(n+1)^{2k}}
 {(3(n-1)^k+(n+1)^k)((n-1)^k+(n+1)^k)}.
\end{align*}
\end{theorem}

\begin{lemma}[Arithmetic bridge]\label{lem:bridge}
For $n\geq2$ and $k\geq1$, with $x,y,z$ as above and
$a=(n-1)^k$, $b=(n+1)^k$, we have
$0<x\leq a<b\leq y\leq z$.
\end{lemma}
\begin{proof}
The standard bounds $0<\varphi(n)<n$ give
$1\leq\varphi(n)\leq n-1$. Each prime factor has positive exponent,
and $p^{e-1}(p+1)=p^e+p^{e-1}>p^e>0$.
There is at least one prime factor when $n\geq2$, so multiplying gives
$\psi(n)>n$, hence $\psi(n)\geq n+1$ by integrality.
The divisor-sum product formula gives
\[
\sigma(n)=\prod_{p\mid n}(1+p+\cdots+p^{e_p}).
\]
Its last two summands dominate the corresponding psi factor;
therefore $\sigma(n)\geq\psi(n)$.
Raising these nonnegative bounds to the positive integer power $k$
preserves weak inequalities and preserves $(n-1)<(n+1)$ strictly.
\end{proof}

\section{Two real inequalities}
Set $r=x/y$, $q=a/b$, and $t=z/y$. Lemma~\ref{lem:bridge} implies
$0<r\leq q<1$ and $t\geq1$: in particular
$x/y\leq a/y\leq a/b$. All denominators below are strictly positive.
Define
\[
 S(r,t)=\frac r{1+t}+\frac1{r+t}+\frac t{r+1}.
\]
The simple expression in Theorem~\ref{thm:simple} equals $S(r,t)$
by division of each numerator and denominator by $y$.
Writing $v=t-1\geq0$, direct common-denominator expansion gives
\begin{align*}
 S(r,t)-S(r,1)
 &=\frac{v\{2v^2+(4+r-r^2)v+r(1-r)(r+3)\}}
 {2(t+1)(r+t)(r+1)},\\
 S(r,1)-S(q,1)
 &=\frac{(q-r)\{4-(1+r)(1+q)\}}{2(1+r)(1+q)}.
\end{align*}
Both differences are nonnegative. Indeed $1-r\geq0$,
$4+r-r^2=4+r(1-r)\geq0$, and
$(1+r)(1+q)\leq4$. Thus
\[
 S(r,t)\geq S(q,1)=\frac{a^2+ab+4b^2}{2b(a+b)}.
\]
Since $ab=(n^2-1)^k$, $a^2=(n-1)^{2k}$ and
$b^2=(n+1)^{2k}$, the last expression is the exact right-hand side of
Theorem~\ref{thm:simple}. This proves that theorem.

For the mixed expression, homogeneity of degree zero gives
\[
 G(r,t)=\frac{r(1+t)}{rt+r+2t}
       +\frac{r+t}{2rt+r+t}
       +\frac{t(r+1)}{rt+2r+t}.
\]
Each mixed numerator and denominator was divided by $y^2$.
Define the polynomial
\begin{align*}
P(r,v)={}&(2r^4+4r^3+7r^2+8r+3)v^2\\
 &+(4r^4-5r^3+8r^2+19r+6)v\\
 &+(1-r)(3r+1)(5r+3).
\end{align*}
The following two identities follow by expansion of common denominators:
\begin{align*}
G(r,t)-G(r,1)
 &=\frac{r(t-1)P(r,t-1)}{D(r,t)},\\
D(r,t)&=(r+1)(3r+1)(rt+r+2t)\\
 &\hspace{15mm}\cdot(rt+2r+t)(2rt+r+t),\\
G(r,1)-G(q,1)
 &=\frac{(q-r)(3+r+q-5rq)}
 {(r+1)(q+1)(3r+1)(3q+1)}.
\end{align*}
The apparent negative coefficient in $P$ is harmless, because
\[
4r^4-5r^3+8r^2+19r+6
=4r^4+5r^2(1-r)+3r^2+19r+6\geq0.
\]
All other factors establishing $P(r,v)\geq0$ are nonnegative.
For the horizontal difference, $rq\leq r$ and $r\leq q\leq1$ give
\[
3+r+q-5rq\geq3-4r+q\geq3-3r\geq0.
\]
Consequently $G(r,t)\geq G(q,1)$, and elementary simplification gives
\[
G(q,1)=\frac{5a^2+5ab+2b^2}{(3a+b)(a+b)}.
\]
The same power identities now give exactly the right-hand side of
Theorem~\ref{thm:mixed}, proving it.

\small
\section{Formal certificate and provenance}
The companion files \texttt{formal/RealRatios.lean} and
\texttt{formal/ArithmeticRatios.lean} prove the identities, signs,
arithmetic bridge, homogeneity, and the two exact final inequalities.
The final declarations are
\texttt{ArithmeticRatios.mathdb\_375256} and
\texttt{ArithmeticRatios.mathdb\_375257}.
They quantify over all natural $n,k$ with $2\leq n$ and $1\leq k$,
interpret all displayed quotients in $\RR$, and use mathlib's actual
\texttt{Nat.totient} and divisor-sum function. Subtraction on the
right-hand sides is real subtraction, so there is no truncated-natural
subtraction in those expressions.

The function \texttt{psi} is the integral product specified above.
The theorem \texttt{psi\_euler\_product} formally proves its equality
with $n\prod_{p\mid n}(1+1/p)$ when $n\ne0$.
The theorem \texttt{sigma\_eq\_sum\_divisors} explicitly identifies
\texttt{sigma} with $\sum_{d\mid n}d$.
The code makes no conjecture-specific assumption and uses no
\texttt{sorry}, newly postulated axiom, or \texttt{native\_decide}.
The replay script runs Lean 4.34.1 with warnings treated as errors;
\texttt{formal/Checks.lean} reports axioms for every local declaration.

Dimitrov's v1 source has the exact statements under labels
\texttt{Conjecture3} and \texttt{Conjecture9}.
Shivam Patel posted both general proofs on 20 August 2026; the exact
post snapshots, identifiers and author fields are in \texttt{sources/}.
Patel's first proof uses derivatives; we replace those steps with explicit
rational certificates. The mixed polynomial is his posted certificate.
These are two distinct source questions sharing an arithmetic bridge,
not a count of their helper lemmas. This is verification of prior posted
proofs, with no claim of a first proof, worldwide priority, or an equality
classification.

\begin{thebibliography}{9}
\bibitem{Dimitrov}
S. I. Dimitrov, \emph{Lower bounds on expressions depending on the
functions $\varphi(n)$, $\psi(n)$ and $\sigma(n)$, III},
arXiv:2606.12484v1, 10 June 2026,
\url{https://arxiv.org/abs/2606.12484v1}.
\bibitem{Patel}
Shivam Patel, posted proofs on MathDB questions 375256 and 375257,
20 August 2026. Exact post snapshots, author fields, timestamps, and
solution identifiers are preserved in the companion \texttt{sources/}
directory.
\end{thebibliography}
\end{document}
